\documentclass[11pt,a4paper]{article}
\usepackage[T1]{fontenc}
\usepackage{lmodern,microtype,amsmath,amssymb,amsthm}
\usepackage[margin=25mm]{geometry}
\usepackage{booktabs,array,xurl}
\usepackage[hidelinks,unicode]{hyperref}
\hypersetup{pdftitle={Erdos 708: Global Dual Normalization and the Remaining Joint Inequality},
pdfauthor={Research working note}}
\setlength{\emergencystretch}{2em}
\newtheorem{theorem}{Theorem}[section]
\newtheorem{lemma}[theorem]{Lemma}
\newtheorem{proposition}[theorem]{Proposition}
\newtheorem{corollary}[theorem]{Corollary}
\theoremstyle{remark}
\newtheorem{remark}[theorem]{Remark}
\DeclareMathOperator{\supp}{supp}
\newcommand{\pos}[1]{(#1)_+}
\newcommand{\LP}{L_{\mathrm{raw}}}
\newcommand{\DD}{\mathcal D}
\title{Erd\H{o}s 708: Global Dual Normalization\\
and the Remaining Joint Inequality}
\author{Research working note}
\date{23 November 2026}
\begin{document}
\maketitle
\begin{abstract}
This note addresses the unrestricted asymptotic target
\(g(n)\le(2+o(1))n\), using the fixed-ambient minimal-obstruction
framework of Section 8 of the preceding report. The target is
\textbf{not proved here}. We prove a uniform normalization of an
optimal raw dual vector: all but \(o(m)\) sources have load above
\(2+2\delta\), their discarded weighted demand is \(o(m)\), and their
entire prime-deletion bank union is saturated in one full optimum.
Its positive weighted row excess over the retained demands is also
\(o(m)\). The proof uses all proper subsets at once and incurs no
per-source rounding multiplier.

We state the exact remaining joint estimate, and prove that even
a weaker version with an additional \(\delta m\) error would close the
general asymptotic problem. We also give a separately verified
counterexample to universal coordinatewise bank-product domination.
That counterexample does not refute the required joint inequality.
No new value of \(g(n)\), restricted source class, or unrestricted
coefficient-two bound is claimed.
\end{abstract}
\tableofcontents

\section{The precise scope and inherited inputs}

For distinct \(A\subseteq\{2,\ldots,M\}\), let \(I\) be a positive
interval of \(M\) consecutive integers, and put
\[
 \tau(A,I)=\min\left\{|B|:B\subseteq I,\
          \prod_{a\in A}a\mid\prod_{b\in B}b\right\}.
\]
The ambient \(M\) and the interval \(I\) are retained when sources are
deleted. In the original definition one may take \(M=\max A\);
the enlarged ambient formulation is needed for induction. A source
equal to one can be deleted without affecting divisibility. Intervals
containing zero are immediate, and a negative interval not containing
zero reduces by absolute values to a positive interval of the same
length.

Write \(P\) for the source primes and
\[
 D_p(A)=\sum_{a\in A}v_p(a),\qquad
 X_w(t)=\sum_{p\in P}w_pv_p(t),\qquad
 H_I(w)=\sum_{b\in I}\pos{X_w(b)-1}.
\]
All valuations are raw, with their full multiplicities. In particular,
target powers are not truncated at \(M\) or at the largest exponent
of a single source.

We use two established inputs of the preceding report \cite{Previous}:
the forced-bank reduction below, and the quadratic ambient theorem
\[
 M\ge16|A|^2\quad\Longrightarrow\quad \tau(A,I)\le2|A|.
 \tag{Q}
\]
The full proof of (Q) was checked in the retained source of Wang
\cite{Wang}; it only uses \(a\le M\), and so applies with the ambient
length fixed after deletion. Its role here is to give
\(M<16m^2\) at a minimal obstruction, not to solve the remaining range.

Fix \(\delta>0\), put \(c=2+3\delta\), and suppose that \(A\) is a
least-cardinality violation at these fixed \(M,I\):
\begin{equation}
 |A|=m,\qquad \tau(A,I)>cm,\qquad
 \tau(S,I)\le c|S|\quad(S\subsetneq A).
 \label{eq:minimal}
\end{equation}
Set \(K=\lfloor c(m-1)\rfloor+1\). If \(q\mid a\) is prime, every
integer deletion cover of \(A\setminus\{a\}\) of size at most \(K-1\)
contains
\[
 C_q(a)=\{b\in I:a/q\mid b\}.
\]
Otherwise add a fresh \(a/q\)-multiple and at most one further target
to repair a missing \(q\)-unit. Full supply guarantees that repair:
\begin{equation}
 \prod A\mid M!\mid\prod I,\qquad
 \prod I=M!\binom{x+M}{M}\quad(I=\{x+1,\ldots,x+M\}).
 \label{eq:supply}
\end{equation}
The resulting number of targets is within the budget \(cm\).
Because \(|C_q(a)|\ge q\), we obtain
\begin{equation}
 p<K<cm\quad(p\in P),\qquad
 r:=|P|\le\pi(cm)=O_\delta(m/\log m).
 \label{eq:primecap}
\end{equation}

The pair-bank inequality from the preceding report states that if
\(pq\mid a\), counting occurrences and allowing \(p=q\), and
\(pq\ge K\), then
\[
 D_p(A)>v_p\!\left(\prod C_q(a)\right)\ge qv_p(a/q).
\]
Consequently the demand core and its complement
\begin{equation}
 E=\{p\in P:p\le\sqrt K\}\cup\{p\in P:D_p(A)>p\},
 \qquad R=P\setminus E
 \label{eq:core}
\end{equation}
have the following properties: every source contains at most one
occurrence from \(R\), that occurrence is the unique largest prime,
and
\begin{equation}
 |E|=O_\delta\!\left(\sqrt{\frac m{\log m}}\right).
 \label{eq:corebound}
\end{equation}
For completeness, the heavy-demand part follows from
\[
 \sum_{D_p(A)>p}p\log p
 <\sum_pD_p(A)\log p
 =\sum_{a\in A}\log a\le m\log M.
\]
Together with \(M<16m^2\) and the elementary prime-count bounds, this
gives \eqref{eq:corebound}. We do not replace that growing core by a
bounded or polylogarithmic core.

\section{The raw LP and the single rounding loss}

The boxed raw relaxation is
\begin{equation}
 \LP(A,I)=\min\left\{\sum_{b\in I}z_b:
 0\le z_b\le1,\quad
 \sum_{b\in I}z_bv_p(b)\ge D_p(A)\ (p\in P)\right\}.
 \label{eq:primal}
\end{equation}
It is feasible by \eqref{eq:supply}. Finite linear programming duality
gives
\begin{equation}
 \LP(A,I)=
 \max_{0\le w_p\le1}\Phi_A(w),\qquad
 \Phi_A(w)=\sum_{a\in A}X_w(a)-H_I(w).
 \label{eq:dual}
\end{equation}
To justify the clipping, when \(w_p>1\), every target with positive
\(p\)-valuation is in the linear part of the hinge. Lowering \(w_p\)
to one cannot decrease the objective because
\(D_p(A)\le\sum_Iv_p\). Thus an optimum exists in the displayed compact
box.

An optimal primal vertex has at most \(r\) strictly fractional
coordinates: otherwise the corresponding columns have a nonzero
linear dependence in the \(r\) rows, giving a feasible two-sided
perturbation. Rounding up yields
\[
 \tau(A,I)\le\LP(A,I)+r.
\]
In the rest of the main argument, fix any optimal clipped dual \(w\),
write \(X=X_w\), and let \(z\) be a primal optimum. Then
\begin{equation}
 F:=\LP(A,I)=\Phi_A(w)>cm-r.
 \label{eq:F-lower}
\end{equation}
The loss \(r=o(m)\) is used once. It will not be added separately for
each deletion.

\section{A uniform normalization using all proper subsets}

\begin{lemma}[Total source shortfall]
\label{lem:shortfall}
Under \eqref{eq:minimal}, for every nonnegative weight vector \(w\),
not necessarily optimal,
\begin{equation}
 \sum_{a\in A}\pos{c-X_w(a)}
 \le\pos{cm-\Phi_A(w)}.
 \label{eq:shortfall}
\end{equation}
In particular, for an optimal clipped vector,
\begin{equation}
 \sum_{a\in A}\pos{c-X(a)}<r.
 \label{eq:optimal-shortfall}
\end{equation}
\end{lemma}

\begin{proof}
Let \(S=\{a\in A:X_w(a)>c\}\). If \(S=A\), the left side of
\eqref{eq:shortfall} is zero. Otherwise \(S\) is a proper subset,
possibly empty. Keep the same weight vector and the same full target
hinge. Weak duality and \eqref{eq:minimal} give
\[
 \Phi_S(w)\le\LP(S,I)\le\tau(S,I)\le c|S|.
\]
Coordinates having zero demand on \(S\) may be retained in this
weak dual objective; their target contributions do not invalidate
weak duality. When \(S=\varnothing\), this reads \(-H_I(w)\le0\).
Now
\[
 \Phi_S(w)-c|S|
 =\Phi_A(w)-cm+\sum_{a\in A\setminus S}(c-X_w(a)).
\]
The last sum is the positive-part shortfall. Its nonnegativity also
shows that \(cm-\Phi_A(w)\ge0\) in this case. This proves
\eqref{eq:shortfall}. Equation \eqref{eq:F-lower} gives
\eqref{eq:optimal-shortfall}, including the case \(F>cm\).
\end{proof}

Put
\begin{equation}
 h=c-\delta=2+2\delta,\qquad
 G=\{a\in A:X(a)>h\},\qquad B=A\setminus G,\qquad
 W=\sum_{a\in B}X(a).
 \label{eq:G}
\end{equation}
Every discarded source contributes at least \(\delta\) to
\eqref{eq:optimal-shortfall}. Therefore
\begin{equation}
 |B|\le\frac r\delta,\qquad
 W\le\frac{(2+2\delta)r}{\delta}=o_\delta(m),\qquad
 |G|=m-O_\delta(r).
 \label{eq:discard}
\end{equation}
Moreover,
\begin{equation}
 \Phi_G(w)>c|G|-r,
 \label{eq:G-lower}
\end{equation}
since
\[
 \Phi_G(w)-c|G|
 =F-cm+\sum_{a\in B}(c-X(a))>-r.
\]
The set \(G\) is a subset used in the dual calculation. We do not
declare it a new integer-minimal instance, and do not assume that
\(w\) is optimal for \(G\). The demand core \eqref{eq:core} is still
defined using the original demands of \(A\).

\section{One global saturated union and its actual demands}

\subsection{The exact optimality identities}

Let
\[
 H=\{b\in I:X(b)>1\},\qquad T=\{b\in I:X(b)=1\}.
\]
Complementary slackness gives
\begin{equation}
 z_b=1\ (b\in H),\quad z_b=0\ (X(b)<1),\quad
 \sum_{b\in I}z_bv_p(b)=D_p(A)\quad(w_p>0).
 \label{eq:cs}
\end{equation}
Consequently
\begin{equation}
 \sum_{a\in A}X(a)=\sum_{b\in H}X(b)+\sum_{b\in T}z_b,
 \qquad F=|H|+\sum_{b\in T}z_b.
 \label{eq:boundary}
\end{equation}
The threshold-one term has not disappeared. A vertex has at most
\(r\) fractional coordinates, but it can have many integral
coordinates with \(X(b)=1\).

For example, take \(A=\{2,4,\ldots,2m\}\), \(M=2m\),
\(I=\{1,\ldots,2m\}\), and \(w_2=1\), with all other weights zero.
Every even target must be selected fully to meet the total
\(2\)-valuation. The optimum has no fractional coordinate and only
one active dual weight, whereas
\(\sum_Tz_b=\lceil m/2\rceil\). The full source-prime count is
\(r=\pi(m)=o(m)\) for \(m\ge2\). This is an exact counterexample to
bounding all threshold-one mass by the number of fractional
coordinates. It is not an instance violating the desired cover bound.

\subsection{Simultaneous bank saturation}

Define the retained union, with all overlaps identified,
\begin{equation}
 U=U(G,I)=\bigcup_{a\in G}\ \bigcup_{\substack{q\mid a\\q\ {\rm prime}}}
                  C_q(a),\qquad
 V_p(U)=\sum_{b\in U}v_p(b).
 \label{eq:union}
\end{equation}
For every \(b\in C_q(a)\),
\[
 X(b)\ge X(a/q)=X(a)-w_q>1+2\delta.
\]
Thus \(U\subseteq H\), and \emph{one full primal optimum} selects
every position of \(U\) with mass one. In particular,
\begin{equation}
 V_p(U)\le D_p(A)\qquad(w_p>0).
 \label{eq:global-rows}
\end{equation}
More generally, any target multiple of a divisor \(d\) satisfying
\(X(d)>1\) is saturated. The prime-deletion union is a useful
explicit subunion, not the whole set of saturated targets.

\begin{lemma}[Joint weighted excess]
\label{lem:excess}
With the original optimal weights, put
\begin{equation}
 \Xi=\sum_{p\in P}w_p\pos{V_p(U)-D_p(G)}.
 \label{eq:Xi}
\end{equation}
Then, uniformly in all the parameters,
\begin{equation}
 \Xi\le W\le\frac{(2+2\delta)r}{\delta}=o_\delta(m).
 \label{eq:Xi-small}
\end{equation}
\end{lemma}

\begin{proof}
For every active row, \eqref{eq:global-rows} gives
\[
 \pos{V_p(U)-D_p(G)}
 \le D_p(A)-D_p(G)=D_p(B).
\]
Multiply by \(w_p\) and sum. The right side is exactly
\(\sum_{a\in B}X(a)=W\). Inactive rows have zero weight.
\end{proof}

This is a sum over actual demands, with no multiplier \(|G|\),
\(|E|\), or \(|R|\). It does not assert that each bank separately
has the full interval budget.

\subsection{Prime powers must be summed before taking row excess}

For a finite set \(S\), write
\(N_S(p^j)=|\{s\in S:p^j\mid s\}|\). The exact level identities are
\[
 D_p(G)=\sum_{j\ge1}N_G(p^j),\qquad
 V_p(U)=\sum_{j\ge1}N_U(p^j).
\]
Hence the excess in \eqref{eq:Xi} is
\begin{equation}
 \Xi=\sum_p w_p
 \left(\sum_{j\ge1}\bigl[N_U(p^j)-N_G(p^j)\bigr]\right)_+.
 \label{eq:level-excess}
\end{equation}
It is not the sum of the positive parts of the separate level
differences. For the elementary valuation example
\(G=\{2,4\}\), \(U=\{8\}\), the differences at powers of two are
\(-1,0,1\). The row excess is zero, whereas the sum of the
positive level differences is one. This example illustrates the
algebraic distinction only; these two sets are not asserted to be
the bank union of a minimal obstruction.

In \eqref{eq:level-excess}, target levels extend as far as the actual
targets require. An unbounded interval shift does not permit their
truncation at the source maximum \(M\). The row balance
\eqref{eq:global-rows} is what controls their selected aggregate.

\subsection{A factorial refinement of an individual bank}

\begin{proposition}
\label{prop:factorial}
Fix \(a\in G\), a prime \(q\mid a\), and an active prime \(p\).
Put \(d=a/q\), \(N=\lfloor M/d\rfloor\). Then
\begin{equation}
 D_p(A)\ge N v_p(d)+v_p(N!).
 \label{eq:factorial}
\end{equation}
In particular, if \(p\mid a/q\) and \(q\ge p\), then
\begin{equation}
 D_p(A)\ge qv_p(a/q)+1>p.
 \label{eq:factorial-cor}
\end{equation}
\end{proposition}

\begin{proof}
The entire \(d\)-bank belongs to \(U\), hence is selected fully.
Its quotients \(b/d\) form \(N'\ge N\) consecutive positive integers.
For every \(j\), their \(p^j\)-multiple count is at least
\(\lfloor N'/p^j\rfloor\). Their total \(p\)-valuation is therefore
at least \(v_p(N'!)\ge v_p(N!)\). Adding the \(N'\) copies of
\(v_p(d)\), then using \eqref{eq:global-rows}, proves
\eqref{eq:factorial}. Since \(N\ge q\ge p\),
\(v_p(N!)\ge1\), proving \eqref{eq:factorial-cor}.
\end{proof}

As a consequence, every active prime \(p\) with \(D_p(A)\le p\)
occurs at most once in a retained source and is its unique largest
prime, even in comparison with inactive prime occurrences. If a
second occurrence \(q\ge p\) existed, including \(q=p\), then
\eqref{eq:factorial-cor} would contradict \(D_p(A)\le p\).
This conclusion does not need the pair threshold \(pq\ge K\).
It still leaves the joint contribution of a potentially large
heavy-demand core unresolved.

\section{The exact remaining step and the closing contradiction}

The sufficient estimate requested as formula (91) in the preceding
report was
\[
 \sup_{0\le w\le1}\left\{
 \sum_{a\in A}\pos{X_w(a)-2-\delta}-H_I(w)\right\}
 =o_\delta(m)
\]
uniformly on the full integer-minimal class. It is not proved here.
The preceding lemmas reduce the obligation to a smaller family of
dual vectors without shrinking the source class.

\subsection{A balanced-optimal sufficient lemma}

Set \(g=|G|\). The following statement is the precise remaining
arithmetic obligation; it is \textbf{unproved}.

\medskip
\noindent\textbf{Joint estimate \(J_\delta\).}
For the original minimal instances and their normalized optimal
data \(w,G,U,\Xi\) above, there exist a constant \(C_\delta\) and
a function \(\varepsilon_\delta(m)\to0\), independent of
\(M,A,x,w\), such that
\begin{equation}
 \boxed{\quad
 \sum_{a\in G}\bigl(X(a)-2-\delta\bigr)
 \le H_I(w)+C_\delta\Xi+
       \varepsilon_\delta(m)m.
 \quad}
 \label{eq:joint}
\end{equation}
The positive part on the source terms is unnecessary because
\(X(a)>2+2\delta\) on \(G\).

\begin{proposition}[Conditional completion]
\label{prop:completion}
If \eqref{eq:joint} holds uniformly, the requested raw bound
\(\LP(A,I)\le(2+\delta)m+o_\delta(m)\) holds on the original
minimal class, giving the unrestricted asymptotic conclusion.
Even the weaker statement obtained by adding \(\delta g\) to the
right side of \eqref{eq:joint} suffices for
\(g(n)\le(2+o(1))n\).
\end{proposition}

\begin{proof}
By \eqref{eq:Xi-small}, the extra \(\Xi\) term is \(o_\delta(m)\).
Equation \eqref{eq:joint} yields
\[
 \Phi_G(w)\le(2+\delta)g+o_\delta(m).
\]
Restore the discarded weighted demand once:
\[
 F=\Phi_G(w)+W
 \le(2+\delta)m+o_\delta(m)
\]
by \eqref{eq:discard}. Rounding costs another \(r=o(m)\), so
\(\tau(A,I)\le(2+\delta)m+o_\delta(m)<(2+3\delta)m\)
for sufficiently large \(m\), contradicting \eqref{eq:minimal}.

For the weaker version, exactly the same argument gives
\[
 F\le(2+2\delta)m+o_\delta(m),
\]
which still leaves a fixed \(\delta m\) gap to the integer budget.
Equivalently, it contradicts \eqref{eq:G-lower} directly.

To pass from minimal instances to the original asymptotic
statement, suppose violations of the \((2+3\delta)n\) budget
exist for arbitrarily large \(n\). Their ambient lengths satisfy
\(M\ge n+1\to\infty\). At each fixed ambient \(M,I\), choose
\eqref{eq:minimal}. The quadratic gate gives
\(m>\sqrt M/4\to\infty\), so the contradiction applies.
Since \(\delta>0\) is arbitrary, the general asymptotic assertion
would follow.
\end{proof}

The weaker statement has three advantages over an all-vector
version of (91): it is needed only at optimal vectors, it has a
single already controlled weighted row-excess allowance, and it
may lose a fixed \(\delta g\). It still requires a new uniform
arithmetic estimate. None of the proved normalization identities
establishes \eqref{eq:joint}.

\subsection{An exact joint accounting identity}

Since \(U\subseteq H\), the retained dual objective can be written
without any repeated target capacity:
\begin{equation}
\begin{split}
 \Phi_G(w)
  ={}& |U|+\sum_p w_p\bigl(D_p(G)-V_p(U)\bigr)\\
    &-\sum_{b\in I\setminus U}\pos{X(b)-1}.
\end{split}
\label{eq:residual}
\end{equation}
Replacing the signed residual sum in this formula by its
coordinatewise positive part increases the right side by exactly
\(\Xi=o_\delta(m)\).

The missing step is a uniform upper bound on
\eqref{eq:residual} by \((2+\delta)g+o(m)\), or by
\((2+2\delta)g+o(m)\) for the weaker completion. The known bank
inequalities give lower bounds on many separate \(V_p(C_q(a))\);
they do not supply this joint bound on \(U\) and its residual
demands. In particular, \(\Xi=o(m)\) controls excess supply,
not the still-positive residual demand.

For arbitrary weights there is also the exact low-load-credit
identity
\begin{equation}
 \left[\sum_{a\in A}\pos{X(a)-t}-H_I(w)\right]
 -\sum_{a\in A}\pos{t-X(a)}
 =\Phi_A(w)-t|A|.
 \label{eq:credit}
\end{equation}
Thus an all-vector hinge bound is stronger than necessary for the
raw optimum. At the optimal vectors under consideration and
\(t\le c\), Lemma~\ref{lem:shortfall} bounds that credit by \(r\).

\section{Why the two tempting aggregation shortcuts fail}

\subsection{Individual banks do not have separate copies of the hinge}

If \(\mu(b)\) is the number of indexed banks \(C_q(a)\) containing
\(b\), then
\[
 \sum_{a\in G}\sum_{\substack{q\mid a\\q\ {\rm prime}}}
       \sum_{b\in C_q(a)}
       \pos{X(b)-1}
 =\sum_{b\in U}\mu(b)\pos{X(b)-1}.
\]
The usable budget is only
\(\sum_{b\in U}\pos{X(b)-1}\). A common multiple can lie in many
banks. Replacing the expression with multiplicity by the single
budget is not a legitimate inequality. Equations
\eqref{eq:Xi} and \eqref{eq:residual} avoid this error, but do not
by themselves prove the missing upper bound.

\subsection{Actual source levels are few, but absolute variation is not quadratic}

Let
\[
 J_E=\sum_{p\in E}\max_{a\in A}v_p(a).
\]
This counts the actual input prime-power levels in the core. In
the present minimal class,
\begin{equation}
 J_E=O_\delta\!\left(\sqrt{\frac m{\log m}}\right),\qquad
 J_E^2+|R|=O_\delta(m/\log m)=o_\delta(m).
 \label{eq:levels}
\end{equation}
Indeed, primes \(p\le m^{1/4}\) contribute at most
\(\pi(m^{1/4})\log_2 M=O(m^{1/4})\). Every larger prime
contributes at most nine levels for sufficiently large \(m\).
Use \eqref{eq:corebound} for their number.

A genuinely joint error \(O_\delta(J_E^2+|R|)\) would therefore
be sufficient. It has not been proved. It cannot be obtained
by merely taking absolute values in a finite divisor expansion:
for \(k\) unit coordinates \(z_i\in\{0,1\}\),
\[
 {\bf1}_{\sum_i z_i>0}
 =\sum_{\varnothing\ne S\subseteq\{1,\ldots,k\}}
       (-1)^{|S|+1}\prod_{i\in S}z_i.
\]
The exact coefficient \(\ell^1\)-norm is \(2^k-1\), although
there are only \(J=k\) levels. This disproves that algebraic
shortcut. It does not disprove a joint estimate exploiting
optimality, actual source demand, and cancellation.

\section{A verified counterexample to coordinatewise bank domination}

Another possible shortcut would be the universal assertion
\begin{equation}
 \prod_{a\in A}a\mid\prod_{b\in U(A,I)}b,
 \qquad
 U(A,I)=\bigcup_{a\in A}
       \bigcup_{\substack{q\mid a\\q\ {\rm prime}}}C_q(a).
 \label{eq:false-bank}
\end{equation}
It is false. The accompanying certificate gives an exact witness;
the following description explains its arithmetic and its limits.

Set
\[
 N=11763=3^2\cdot1307,\quad H_0=5881,\quad p=11777,
 \qquad S=\{s\le N:s\text{ is composite}\}.
\]
Then \(\pi(N)=1409\) and \(|S|=10353\). The file
\nolinkurl{bank_product_certificate.json} specifies CRT residues
\(c_r\) for every prime \(r\le N\), and a center \(C\) satisfying
\[
 C\equiv c_r\pmod r,\qquad C\equiv H_0+1\pmod p.
\]
It gives 1411 offsets in \([-H_0,H_0]\) avoiding every prime
at most \(N\). The independent exact verifier reconstructs
the CRT solution and checks all these residues. For
\(J=\{C-H_0,\ldots,C+H_0\}\), the proper prime-deletion
neighborhood of \(S\) is
\[
 V=\{t\in J:\text{some prime }r\le N/2\text{ divides }t\},
 \qquad |V|=10352.
\]
To identify this neighborhood, every quotient \(s/q\) is at
least two and has a prime factor at most \(N/2\). Conversely
each prime \(r\le N/2\) is the quotient of the source \(2r\)
by two.

Lift to
\[
 A=pS,\qquad M=pN=138532851,\qquad
 I=\{p(C-H_0-1)+1,\ldots,p(C+H_0)\}.
\]
The \(p\)-multiples in \(I\) are precisely \(pt\), \(t\in J\).
The last CRT congruence makes the residues of \(J\) modulo \(p\)
equal to \(1,\ldots,N\), so every such target has \(p\)-valuation
one. Deleting a prime \(q\mid s\) from \(ps\) gives the
neighborhood condition \(s/q\mid t\). Deleting \(p\) gives
\(s\mid t\), which adds no new \(p\)-multiple. Therefore
\[
 U(A,I)\cap p\mathbb Z=\{pt:t\in V\},
\]
and exactly
\begin{equation}
 D_p(A)=10353,\qquad V_p(U(A,I))=10352.
 \label{eq:counter}
\end{equation}
Every source has at least three prime occurrences.

\paragraph{The scope is essential.}
This family is not a minimal violation. The source \(4p\) occurs,
and its bank obtained by deleting \(p\) contains all multiples
of four, at least
\[
 \lfloor M/4\rfloor=34633212>10|A|.
 \]
It cannot fit inside the relevant deletion cover. The example
even has \(p<2|A|-1\) and \(M<16|A|^2\), illustrating why those
two numerical restrictions alone do not replace full minimality.

It also does not refute the joint weighted expansion at high
source loads. For every other source prime \(r\ne p\), the bank
of four gives
\[
 V_r(U)>1471D_r(A).
\]
Indeed \(D_r(A)\le N/(r-1)\), whereas
\[
 V_r(U)\ge\left\lfloor\frac{pN}{4r}\right\rfloor
 \ge\frac p8D_r(A)-1
 \ge\left(\frac p8-1\right)D_r(A)>1471D_r(A).
\]
If clipped weights satisfy \(X(a)>2\) for every source, then
\(\sum_{r\ne p}w_rD_r(A)>(2-w_p)|A|\ge|A|\).
The weighted total bank surplus is consequently positive.
The negative \(p\)-row in \eqref{eq:counter} cannot be used as a
counterexample to the required weighted statement.

\paragraph{Verification record.}
The standalone standard-library verifier checked 1410 CRT
conditions, 1,988,099 rough-residue avoidance conditions, and
27,390 prime-deletion quotients. It returns the exact deficit
one in \eqref{eq:counter}. The certificate SHA256 is
\begin{center}
\small\texttt{6da6d75f3f7a2580c0eb0cb6ac1c4d2c17c4358b2f9ac1ae82a66941956ca92b}.
\end{center}
Only the quotient interval of length 11763 is enumerated.
No numerical search for a new value of \(g(n)\) is involved.

\section{Status and uniformity audit}

The following statements are proved in this note on the full
minimal class:
\begin{itemize}
 \item the total optimal source shortfall is less than the single
 rounding loss \(r=o(m)\);
 \item all but \(o(m)\) sources have load above \(2+2\delta\),
 and their discarded weighted demand is \(o(m)\);
 \item their one global prime-deletion bank union is fully
 saturated in a full optimal primal solution;
 \item its positive weighted row excess over the retained
 demands is \(o(m)\);
 \item each active bank obeys the factorial lower bound
 \eqref{eq:factorial}, uniformly in the interval shift;
 \item the actual source core-level count satisfies
 \eqref{eq:levels}.
\end{itemize}
The arithmetic estimate \eqref{eq:joint}, its weaker version
with \(\delta g\), and a sufficient weighted bank-expansion
replacement remain unproved. Formula (91) is neither proved
nor refuted. The exact place where the current argument stops
is the bound on the joint residual functional
\eqref{eq:residual}.

No small-core hypothesis, bound on the number of substantial
weights, or gap between prime sizes is imposed in the proved
reductions. The interval length and shift are not altered.
All error constants may depend on fixed \(\delta\), but not
on \(M,A,x,w\). The target prime powers are retained in full.
The only use of an asymptotic estimate is through bounds
depending on \(m\) alone, including \(r=O_\delta(m/\log m)\).

Even a proof of the remaining joint estimate would establish
the asymptotic coefficient two, not automatically the exact
integer inequality \(2n\). The latter would still require
removing the asymptotic and rounding slack or an exact
construction. Apart from inherited exact results such as
(Q), no new exact \(2n\) assertion is claimed here.

\appendix
\section{A separately valid raw-minimal alternative}

The main argument deliberately kept the original integer-minimal
instance. There is also a fully justified raw-minimal reduction,
which may be useful for a subsequent attack.

Fix \(c_0>2,M,I\), and take \(A_0\) of least cardinality
\(m_0\) with \(\LP(A_0,I)>c_0m_0\). Write
\[
 F_0=c_0m_0+e,\quad e>0,\qquad K_0=c_0(m_0-1)+1.
\]
Every deletion has a fractional cover of cost at most \(K_0-1\).
For the same optimal full dual, weak duality on each deletion gives
\[
 X(a)\ge c_0+e>c_0\qquad(a\in A_0).
 \tag{A1}
\]

\begin{lemma}[Residual unit repair and union saturation]
For every deletion fractional cover \(y\) of cost at most \(K_0-1\),
and \(U_a=\bigcup_{\substack{q\mid a\\q\ {\rm prime}}}C_q(a)\),
\[
 \sum_{b\in U_a}(1-y_b)<1.
 \tag{A2}
\]
\end{lemma}
\begin{proof}
A partial fractional selection with positive valuation deficits
\(d_p\) can be extended at cost at most \(\sum_pd_p\).
Whenever a row is deficient, full supply supplies enough residual
valuation in the capacities \(1-y_b\). A positive integer
valuation is at least one, so buying the needed valuation costs
no more than the deficit. Sequential repairs never increase
other deficits.

If (A2) failed, buy a total mass one inside \(U_a\). Label each
used target by a witnessing removed prime \(q\).
Its valuations supply \(v(a)-{\bf e}_q\), so the mixed mass
supplies \(v(a)\) except for a nonnegative vector with total
mass at most one. Repair it at additional cost at most one.
The resulting cover costs at most
\(c_0(m_0-1)+2<c_0m_0\), a contradiction.
\end{proof}

Since each \(q\)-bank has at least \(q\) targets, (A2) gives
\(q<K_0<c_0m_0\). If \(pq\mid a\) and \(pq\ge K_0\), then
\[
 D_p(A_0)>(q-1)v_p(a/q).
 \tag{A3}
\]
Indeed, otherwise the almost-full \(a/q\)-bank already supplies
the full \(p\)-demand. The \(a/(pq)\)-bank has residual mass at
least \(pq-(K_0-1)\ge1\). Buy mass one there and repair at most
one \(q\)-unit, again below the failing budget.
The correct factor is \(q-1\) in this fractional argument.

Thus the core
\[
 E_0=\{p\le\sqrt{K_0}\}\cup\{p:D_p(A_0)\ge p\}
\]
has the same size bound as before, and its complement has
at most one prime occurrence per source, as its unique
largest prime. The slight change from \(D_p>p\) to \(D_p\ge p\)
is necessary when using (A3).

For any such deletion cover \(y\), with cost \(C_a\), the exact
primal--dual gap identity is
\begin{align*}
 C_a-F_0+X(a)
 ={}&\sum_p w_p\left(\sum_by_bv_p(b)-D_p(A_0)+v_p(a)\right)\\
 &+\sum_{X(b)>1}(1-y_b)(X(b)-1)
   +\sum_{X(b)<1}y_b(1-X(b)).
 \tag{A4}
\end{align*}
All three terms are nonnegative, and the left side is at most
\(X(a)-c_0-e\). The weighted oversupply and the two types of
target slack share this one budget. If
\(w_{\max}(a)=\max_{\substack{q\mid a\\q\ {\rm prime}}}w_q\),
(A4) sharpens (A2) to
\[
 \sum_{b\in U_a}(1-y_b)
 \le\frac{X(a)-c_0-e}{X(a)-1-w_{\max}(a)}<1.
 \tag{A5}
\]

Finally, this alternative is reached from the requested integer
budget with the same fixed \(\delta\). Start with
\eqref{eq:minimal}. Equation \eqref{eq:F-lower} gives
 \(F>(2+2\delta)m\) for large \(m\).
At the same \(M,I\), choose a smallest subset with raw value
above \((2+2\delta)\) times its cardinality. Apply (A1)--(A5)
with \(c_0=2+2\delta\). Its own quadratic gate gives
\(m_0>\sqrt M/4\to\infty\), and its own prime cap uses \(m_0\),
not the size of the original source set.
A uniform upper bound \((2+\delta)m_0+o(m_0)\) would contradict
this raw minimality. This is an alternative formulation of the
remaining general task, not a proof of that upper bound.

\begin{thebibliography}{9}
\bibitem{Previous}
Research working report,
\emph{Erd\H{o}s 708: continuation report}, 23 November 2026,
Section 8, equation (91). Retained as
\nolinkurl{erdos708_continuation.pdf} in the preceding deliverable.
\bibitem{Wang}
Ziming Wang,
\emph{Fractional-to-Integer Transference and Five-Prime Bounds
for Interval Divisor Covers}, appendix on the exact quadratic gate.
Primary project and result catalogue:
\url{https://github.com/science-prince/erdos-708-interval-divisor-covers}
and
\url{https://github.com/science-prince/erdos-708-interval-divisor-covers/blob/main/docs/RESULTS.md}.
The retained proof, rather than a new status claim, is the input
used in this note.
\end{thebibliography}
\end{document}
